% Generated draft: requires mathematical and source review.
\documentclass[11pt]{article}
\usepackage[margin=25mm]{geometry}
\usepackage{fontspec}
\setmainfont{Latin Modern Roman}
\usepackage{amsmath,amssymb,mathtools,unicode-math}
\setmathfont{Latin Modern Math}
\usepackage{enumitem,needspace,xurl,hyperref}
\setlength{\parindent}{0pt}
\setlength{\parskip}{5pt}
\setlength{\emergencystretch}{3em}
\newcommand{\N}{\mathbb{N}}
\newcommand{\Z}{\mathbb{Z}}
\newcommand{\Q}{\mathbb{Q}}
\newcommand{\R}{\mathbb{R}}
\newcommand{\C}{\mathbb{C}}
\newcommand{\F}{\mathbb{F}}
\newcommand{\A}{\mathbb{A}}
\newcommand{\PP}{\mathbb{P}}
\newcommand{\E}{\mathbb{E}}
\newcommand{\eps}{\varepsilon}
\newcommand{\dd}{\,\mathrm{d}}
\begin{document}
\section*{UnsolvedMath ID: 30000494}
\textbf{Source problem code:} catalogue entry
\subsection{Definitions and mathematical statement}
Let $p$ be a prime and let $K$ be a complete discretely valued field whose residue field has characteristic $p$. Fix a separable closure $K^{\mathrm{sep}}$, let $K(p)$ be the compositum of the finite Galois extensions of $K$ of $p$-power degree, and define the profinite group
\[
 \Gamma_K(p)=\operatorname{Gal}(K(p)/K)
\]
with its Krull topology. This is the maximal pro-$p$ Galois group; it is not an assertion that the whole absolute Galois group is pro-$p$.

Determine the structure of $\Gamma_K(p)$, especially when $K$ is an $N$-dimensional local field. A usual higher-local-field convention gives a chain $K=K_N,K_{N-1},\ldots,K_0$ such that each $K_i$ for $i\ge1$ is a complete discretely valued field with residue field $K_{i-1}$, and $K_0$ is finite. The report's exact convention must be checked. The original hypothesis that the residue field of $K$ has characteristic $p$ remains in force; it is not erased merely by specifying a higher-dimensional local field.

The requested structure concerns the entire topological group. A description might take the form of a presentation or an intrinsic characterization in terms of the field data, but no required form or conjectural classification is supplied. The record does not restrict to equal characteristic, mixed characteristic, perfect residue fields, fields containing roots of unity, or finitely generated pro-$p$ groups. No such hypotheses or finiteness conclusion are added.

The supplied assessment reports explicit Lie-theoretic descriptions and presentations for substantial quotients of period $p$ and nilpotence class less than $p$. Such quotients are partial information and are not substituted for the full group. Their precise construction, parameter restrictions, and whether additional ramification information is part of the desired full structure are not included.
\subsection{Short English statement}
The question asks for the full maximal pro-$p$ Galois group of a complete discretely valued field with residue characteristic $p$, particularly in higher local dimension. The supplied quotient-level Lie descriptions are retained as partial context rather than a classification of the entire group.
\subsection{Sources}
\begin{itemize}
\item Local UnsolvedMath registry, permanent ID 30000494, code catalogue entry, ``Maximal Pro-$p$ Galois Groups of Higher Local Fields.'' Proposed year 2006; proposer unspecified.
\item \emph{Pro-$p$ Extensions of Global Fields and pro-$p$ Groups} (2007), Oberwolfach Reports, pp.\ 1463--1536, DOI 10.4171/owr/2006/25; \url{https://doi.org/10.4171/owr/2006/25}. The displayed report year 2007 differs from the DOI and proposed year 2006; escaped dollar signs in the title are normalized as mathematical $p$.
\item Supplied ``Abrashkin 2024 primary paper'': \url{https://doi.org/10.5802/jtnb.1293}.
\item Supplied ``Abrashkin higher-local-field results'': \url{https://arxiv.org/abs/2009.02501}.
\item Statement status corrected\_verified; the supplied verification says editorial and LaTeX cleanup preserved meaning and added no source-dependent content.
\item Literature triage checked 2026-08-21: partially solved. It says the full higher-local-field maximal pro-$p$ group is not classified while important period-$p$, class-less-than-$p$ quotients have explicit descriptions. Exact scope and current status require independent verification. Created 2026-08-25; updated 2026-09-16.
\item Registry curation by Ulam AI / UnsolvedMath Contributors, CC BY 4.0; underlying report and literature provenance remain as supplied. Unreviewed statement-only draft.
\end{itemize}
\end{document}
